% Warmup 必要性：横轴 epoch（前5 epoch warmup），纵轴 loss
% 对比"无 warmup"（早期发散/震荡）vs "有 warmup"（平滑下降）
\documentclass[tikz,border=4pt]{standalone}
\usepackage{ctex}
\usepackage{amsmath}
\usetikzlibrary{calc, arrows.meta, decorations.pathmorphing}

\begin{document}
\begin{tikzpicture}[scale=1.0, thick]

% 坐标轴
% 横轴：epoch 0..50，映射到 x=0..8.5
% 纵轴：loss 0..5，映射到 y=0..4.5
\draw[->, gray] (-0.2, 0) -- (9.0, 0) node[right, font=\small] {Epoch};
\draw[->, gray] (0, -0.2) -- (0, 4.8) node[above, font=\small] {训练 Loss};

% 横轴刻度
\foreach \x/\lab in {0/0, 0.85/5, 1.7/10, 4.25/25, 8.5/50}{
  \draw[gray, thin] (\x, 0) -- (\x, -0.07);
  \node[below, font=\footnotesize, gray] at (\x, -0.1) {\lab};
}
% 纵轴刻度
\foreach \y/\lab in {0/0, 1/1, 2/2, 3/3, 4/4, 4.5/5}{
  \draw[gray, thin] (0, \y) -- (-0.07, \y);
  \node[left, font=\footnotesize, gray] at (-0.1, \y) {\lab};
}

% Warmup 区域（前5 epoch = x=0..0.85）阴影
\fill[red!10] (0, 0) rectangle (0.85, 4.8);
\draw[red!40, thin, dashed] (0.85, 0) -- (0.85, 4.8);
\node[above, font=\footnotesize, red!60] at (0.43, 4.6) {Warmup\\(前5epoch)};

% ── 曲线1：无 Warmup（红色，震荡/发散）────────────────────
% 前5 epoch：剧烈震荡
% 模拟：大起伏随机游走后缓慢下降
\draw[red, very thick]
  plot[smooth, tension=0.5]
  coordinates {
    (0.00, 3.5)
    (0.15, 4.2)
    (0.25, 3.0)
    (0.35, 4.5)
    (0.45, 3.2)
    (0.55, 4.0)
    (0.65, 3.5)
    (0.75, 4.3)
    (0.85, 3.8)
    (1.10, 3.2)
    (1.50, 2.6)
    (2.00, 2.1)
    (2.80, 1.7)
    (3.70, 1.4)
    (4.80, 1.2)
    (6.00, 1.05)
    (7.20, 0.95)
    (8.50, 0.88)
  };

% 震荡标注
\draw[<-, red!70, thin] (0.35, 4.5) -- (0.6, 4.7) node[right, font=\footnotesize, red!70] {剧烈震荡};

% ── 曲线2：有 Warmup（蓝色，平滑下降）─────────────────────
% Warmup 阶段（0..0.85）：从高loss平滑快速下降
% 之后正常衰减
\draw[blue, very thick]
  plot[smooth, tension=0.8]
  coordinates {
    (0.00, 3.5)
    (0.15, 3.3)
    (0.30, 3.0)
    (0.50, 2.7)
    (0.70, 2.4)
    (0.85, 2.1)
    (1.10, 1.8)
    (1.50, 1.5)
    (2.00, 1.3)
    (2.80, 1.1)
    (3.70, 0.95)
    (4.80, 0.85)
    (6.00, 0.78)
    (7.20, 0.73)
    (8.50, 0.70)
  };

% 平滑标注
\draw[<-, blue, thin] (1.2, 1.75) -- (2.2, 2.2) node[right, font=\footnotesize, blue] {平滑快速下降};

% 最终 loss 比较
\draw[red, thin, dashed] (8.5, 0.88) -- (9.0, 0.88);
\draw[blue, thin, dashed] (8.5, 0.70) -- (9.0, 0.70);
\draw[<->, black!50, thin] (8.85, 0.70) -- (8.85, 0.88);
\node[right, font=\tiny, black!50] at (9.0, 0.79) {更低};

% ── 图例 ──────────────────────────────────────────────────
\begin{scope}[xshift=2.8cm, yshift=0.2cm]
  \draw[red, very thick] (0, 4.3) -- (1.0, 4.3);
  \node[right, font=\footnotesize, red] at (1.0, 4.3) {无 Warmup（震荡）};

  \draw[blue, very thick] (0, 3.9) -- (1.0, 3.9);
  \node[right, font=\footnotesize, blue] at (1.0, 3.9) {有 Warmup（平滑）};

  \draw[black!30, thin, rounded corners=2pt] (-0.2, 3.7) rectangle (4.2, 4.5);
\end{scope}

% Warmup 机制解释
\node[font=\footnotesize, black!60, align=left] at (5.5, 4.2)
  {Warmup：$\eta_t = \eta_{\max} \cdot \frac{t}{T_w}$};

% ── 标题 ──────────────────────────────────────────────────
\node[font=\small\bfseries] at (4.25, -0.65) {Warmup 的必要性：有/无 Warmup 的训练 Loss 对比};

\end{tikzpicture}
\end{document}
